\scanpage{18}
% Reused typescript, rotated 180 degrees in the scan.
\begin{center}-2.28-\end{center}
\[
L^N=\phi^*(\underline{o}(1))
\]
Therefore:
\[
\mu^L(x,\lambda)=\frac1N\mu^{o(1)}(\phi x,\lambda).
\]
As in \S1, the action of $G$ on $\mathbb P_{n-1}$ is induced by a linear representation of $G$ in the affine cone $\mathbb A^n$ over $\mathbb P_{n-1}$. Suppose we choose coordinates in $\mathbb A^n$ so that the representation of $T$ is diagonalized: i.e.\ $\alpha\in T_k$ acts via the matrix $\{\chi_i(\alpha)\cdot\delta_{ij}\}$ for suitable characters $\chi_1,\ldots,\chi_n$ of $T$. Let $\chi_i$ induce the linear functional $\bar\chi_i$ on $\unsure{\Gamma|\Pi}^{\mathbb R}(T)$, and as usual let a 1-PS $\lambda$ of $T$ induce the point $\bar\lambda\in\unsure{\Gamma|\Pi}^{\mathbb R}(T)$. Then if $(\phi x_1^*,\ldots,\phi x_n^*)$ are coordinates of a homogeneous point lying over $\phi x$, Proposition 2.3 asserts:
\[
\mu^{o(1)}(\phi x,\lambda)=\max\{-\bar\chi_i(\bar\lambda)\mid i\text{ such that }\phi x_i^*\neq0\}.
\]
We conclude:

\uline{Proposition 2.14}: Let a reductive $G$ act on $X$, proper over $k$. Let $T\subset G$ be a maximal torus, and let $L\in\operatorname{Pic}^G(X)$ be ample. Then there is a finite set of linear functionals $\ell_1,\ldots,\ell_n$ on $\unsure{\Gamma|\Pi}^{\mathbb R}(T)$ which are rational on $\unsure{\Gamma|\Pi}^{\mathbb Q}(T)$ with the following property:

\begin{quote}
for every $x\in X_k$, there is a subset $I\subset\{1,\ldots,n\}$ such that $\mu^L(x,\lambda)=\max\{\ell_i(\bar\lambda)\mid i\in I\}$ for all 1-PS's $\lambda$ of $T$.
\end{quote}

\uline{Corollary 2.15}: If $\lambda_1,\lambda_2$ are 1-PS's of $T$ and $x\in X_k$, then
\[
\mu^L(x,\lambda_1\circ\lambda_2)\leq\mu^L(x,\lambda_1)+\mu^L(x,\lambda_2)
\]
\mspage{17}
